\section{INTRODUÇÃO}

\subsection{Contexto e Objetivo do Documento}

O presente documento constitui o Documento de Projeto de Software (SDD) do sistema \textbf{MedAgenda}, desenvolvido como trabalho final (T3) da disciplina Engenharia de Software I, ministrada pelo Prof. Me. Julio Leite Azancort Neto na Universidade Federal do Pará, no semestre 2026.1.

Este documento representa a terceira e última etapa do ciclo de Engenharia de Software praticado ao longo do semestre, correspondendo à fase de \textbf{Projeto de Software}, conforme fundamentado na literatura da área \cite{sommerville2011, pressman2016}. O objetivo principal é demonstrar a capacidade de projetar, modelar e especificar um sistema de software completo, aplicando os conceitos estudados na disciplina, incluindo modelagem UML, projeto arquitetural, princípios de orientação a objetos, estratégias de verificação e validação, e projeto de interface com foco em usabilidade.

\subsection{Identificação do Sistema}

O \textbf{MedAgenda} é um sistema web de agendamento e gestão de consultas médicas voltado para clínicas de pequeno e médio porte. O sistema permite que \textbf{pacientes} realizem agendamentos de consultas de forma online, \textbf{médicos} gerenciem suas agendas e registrem prontuários, e \textbf{administradores} controlem cadastros de profissionais, configurações da clínica e geração de relatórios gerenciais.

A motivação para a escolha deste sistema reside na sua relevância prática, pois a digitalização de processos em clínicas médicas é uma demanda crescente no Brasil, além da clareza de funcionalidades que permite exercitar de forma objetiva todos os conceitos exigidos pelo trabalho: atores bem definidos, objetos com ciclos de vida claros (consultas e pacientes), e interfaces adaptadas para cada perfil de usuário.

\subsection{Justificativa de Escolha do Sistema}

Este trabalho apresenta um sistema novo, não utilizado nas etapas anteriores (T1 e T2). A decisão pela troca de sistema foi motivada pela identificação de limitações no sistema anteriormente analisado, que apresentava um escopo restrito de funcionalidades, dificultando a produção dos artefatos UML e de arquitetura exigidos para o T3. O MedAgenda, por outro lado, oferece um domínio intuitivo e completo, com três perfis de usuário distintos e fluxos essenciais de clínica médica (agendamento, cancelamento, prontuário e relatórios), permitindo a produção de artefatos claros e bem estruturados.
